\section{Introduction}
\label{sec:intro}

Removing transformer blocks and then continuing pretraining (CPT, often called
``healing'') is an attractive route to smaller language models
~\citep{gromov2024unreasonable,shortenedllama,men2024shortgpt,
yang2024laco}. A pruned model can improve rapidly on held-out next-token loss,
and endpoint perplexity is therefore a natural measure of whether retraining is
working. It is not, however, a complete description of what the structural
intervention removed or what CPT restored.

Perplexity averages local prediction loss over a corpus. Factual recall,
multiple-choice answer selection, generation, and task-specific interfaces
probe different observables. Compression studies have documented that
aggregate language-model metrics can miss changes in parametric knowledge,
downstream behavior, or safety~\citep{costcompression,beyondperplexity,
jaiswal2024truth}. Prior depth-pruning studies also report healing curves and
loss--task dissociations~\citep{gromov2024unreasonable,shortenedllama,minitron,
iterabre}. We do not claim either phenomenon as new. Instead, we ask what can be
learned from one densely measured OLMo prune--regrow path when it is read
alongside several imperfect but complementary controls.

\paragraph{Research question.}
We ask: \emph{what does continued pretraining restore in the observed OLMo
prune--regrow runs, and which controls change the interpretation?}
We treat recovery as a multi-axis object---held-out likelihood, target
capability, scoring interface, exact construction, and training budget---rather
than a scalar certified by perplexity.

We instantiate this case study on OLMo-2-1124-7B~\citep{olmo2}. Prefix arms retain
$k$ pretrained decoder blocks, append two randomly initialized blocks, and
continue pretraining on a fixed Dolmino corpus array and nominal schedule. The principal keep14 arm is
tracked to 200k steps. Controls include the intact base, an unpruned model
continued on the same corpus array and nominal schedule to an available 25k
checkpoint, a
frozen-prefix arm, a fully random same-shape model, and a non-contiguous
ShortGPT-16 endpoint. Evaluation combines held-out likelihood, a broad zero-shot
likelihood suite, answer-letter and complete-option MMLU, and no-retrieval
PopQA, TriviaQA, and NQ-open.

The central observation is bounded and recurs across the tested interfaces. In
this run, keep14 reaches
PPL 10.561, yet remains far below the intact model on answer-letter MMLU and
all three closed-book tasks. Complete-option scoring narrows the MMLU gap,
consistent with an answer-symbol or readout component, but the random-init
content-score floor argues against interpreting that protocol as recovered subject
knowledge. ShortGPT is substantially stronger on both PPL and MMLU, so the
keep14 endpoint is not shared by every tested 16-layer construction. Coupled
structural differences prevent a single-factor explanation, and all central
training comparisons are single realizations.

\paragraph{Contributions.}
\begin{itemize}[leftmargin=1.2em,itemsep=2pt]
\item \textbf{An OLMo prune--regrow recovery path.} We jointly track
in-domain likelihood, answer-letter MMLU, and broader behavior over available
checkpoints, with the single-run scope explicit.
\item \textbf{A complementary control bundle.} An available short-horizon
full32 branch, same-shape operating points, two MMLU interfaces, closed-book
generation, and a coupled ShortGPT construction comparison delimit competing
interpretations without claiming clean factor isolation.
\item \textbf{A case-study reporting proposal.} The measurements motivate
reporting likelihood, target capability, interface, exact construction, and
recovery budget as separate axes; they do not validate a universal protocol.
\end{itemize}
